% ==============================================================================
% ================================= Aufgabe 2 =================================
% ==============================================================================

\chapter{Aufgabe 2: Projektstruktur und Datenbankanbindung}
\label{chap:aufgabe2}
\thispagestyle{fancy}

In dieser Aufgabe wird die eingerichtete Projektstruktur für die Umfrageanwendung verwendet und die Verbindung zur Datenbank aus Aufgabe 1 implementiert.

\section{Projektstruktur}

Die folgende Darstellung zeigt den Pfad, der im \texttt{\acs{htdocs}}-Verzeichnis von \acs{XAMPP} eingerichtet wurde:

\begin{forest}
    for tree={
        font=\ttfamily,
        grow'=0,
        child anchor=west,
        parent anchor=south,
        anchor=west,
        calign=first,
        edge path={
            \noexpand\path [draw, \forestoption{edge}] (!u.south west) +(10pt,0) |- node[fill,inner sep=1.5pt] {} (.child anchor)\forestoption{edge label};
        },
        before typesetting nodes={
            if n=1{
                insert before={[,phantom,calign with current]}
            }{}
        },
        fit=band,
        before computing xy={l=22pt},
    }
    [\textbf{C:\textbackslash xampp\textbackslash htdocs\textbackslash survey}
        [\textbf{inc\textbackslash}
            [\textbf{db.php}]
        ]
    ]
\end{forest}

\begin{itemize}
    \item Der Ordner \texttt{survey} dient als Hauptverzeichnis der Anwendung.
    \item Der Unterordner \texttt{\acs{inc}} (\textit{für \textit{include}}) enthält wiederverwendbare \acs{PHP}-Komponenten.
    \item Die Datei \texttt{db.php} implementiert die zentrale Datenbankverbindung.
\end{itemize}

Diese Struktur entspricht dem gängigen \acs{MVC}-ähnlichen Aufbau, bei dem Konfigurations- und Verbindungsdateien zentral in einem \texttt{\acs{inc}}-Ordner verwaltet werden.

\section{Datei \texttt{db.php}}

Die Datei \texttt{db.php} enthält alle erforderlichen Informationen für die Verbindung zur My\acs{SQL}-Datenbank. Sie verwendet die \acs{PHP}-Erweiterung \texttt{PDO} (\textit{\acs{PHP} Data Objects}), die eine sichere und objektorientierte Schnittstelle für den Datenbankzugriff bietet.

\subsection{Vollständiger Code}

Der vollständige Quellcode der Datei \texttt{db.php} lautet:

\lstinputlisting[
    language=php, 
    caption={Datenbankverbindung in \texttt{db.php}}, 
    label={lst:db-php}
]{asset/code/Aufgabe2.tex}

\subsection{Erläuterung der einzelnen Bestandteile}

\subsubsection{Konfigurationsparameter}

Die Verbindungsparameter werden als \acs{PHP}-Variablen definiert:

\begin{table}[h]
\centering
\begin{tabular}{|l|l|l|}
\hline
\textbf{Variable} & \textbf{Wert} & \textbf{Bedeutung} \\
\hline
\texttt{\$host} & \texttt{localhost} & Datenbankserver (lokal bei \acs{XAMPP}) \\
\texttt{\$database} & \texttt{umfrage\_db} & Name der Datenbank aus Aufgabe 1 \\
\texttt{\$username} & \texttt{root} & Standard-Benutzer bei \acs{XAMPP} \\
\texttt{\$password} & \texttt{''} (leer) & Standard-Passwort bei \acs{XAMPP} \\
\hline
\end{tabular}
\caption{Konfigurationsparameter der Datenbankverbindung}
\label{tab:db-config}
\end{table}

\subsubsection{\acs{PDO}-Verbindung}

Die Verbindung wird mit folgendem Aufruf hergestellt:

\begin{lstlisting}[language=php]
$pdo = new PDO("mysql:host=$host;dbname=$database;charset=utf8mb4", $username, $password);
\end{lstlisting}

Der DSN (\textit{Data Source Name}) enthält:
\begin{itemize}
    \item \texttt{mysql:} – Treiber für My\acs{SQL}
    \item \texttt{host=\$host} – Server-Adresse
    \item \texttt{dbname=\$database} – Datenbankname
    \item \texttt{charset=utf8mb4} – Zeichensatz (\textit{unterstützt auch Emojis und Sonderzeichen})
\end{itemize}

\subsubsection{\acs{PDO}-Konfiguration}

Nach der Verbindung werden folgende Attribute gesetzt:

\begin{lstlisting}[language=php]
$pdo->setAttribute(PDO::ATTR_ERRMODE, PDO::ERRMODE_EXCEPTION);
$pdo->setAttribute(PDO::ATTR_DEFAULT_FETCH_MODE, PDO::FETCH_ASSOC);
\end{lstlisting}

\begin{itemize}
    \item \texttt{PDO::ATTR\_ERRMODE} auf \texttt{PDO::ERRMODE\_EXCEPTION}: Fehler werden als Exceptions geworfen, was eine saubere Fehlerbehandlung ermöglicht.
    \item \texttt{PDO::ATTR\_DEFAULT\_FETCH\_MODE} auf \texttt{PDO::FETCH\_ASSOC}: Abfrageergebnisse werden standardmäßig als assoziative Arrays zurückgegeben.
\end{itemize}

\subsubsection{Fehlerbehandlung}

Die Verbindung wird in einem \texttt{try-catch}-Block durchgeführt:

\begin{lstlisting}[language=php]
try {
    // Verbindungsaufbau
} catch (PDOException $e) {
    die("Datenbankverbindung fehlgeschlagen: " . $e->getMessage());
}
\end{lstlisting}

Bei einem Verbindungsfehler wird eine aussagekräftige Fehlermeldung ausgegeben und die Skriptausführung abgebrochen. Dies verhindert, dass die Anwendung mit einer defekten Datenbankverbindung weiterläuft.

\subsection{Einbindung über \texttt{require}}

Die Datei \texttt{db.php} wird in anderen \acs{PHP}-Skripten über die \texttt{require}-Anweisung eingebunden:

\begin{lstlisting}[language=php]
require __DIR__ . '/inc/db.php';
\end{lstlisting}

\texttt{require} wird verwendet, da es im Fehlerfall eine \texttt{E\_COMPILE\_ERROR} auslöst und die Skriptausführung abbricht – im Gegensatz zu \texttt{include}, das nur eine Warnung generiert. Dies ist für die Datenbankverbindung kritisch, da die Anwendung ohne funktionierende Verbindung nicht nutzbar wäre.

\section{Test der Datenbankverbindung}

Um die Funktionalität der Verbindung zu überprüfen, wurde eine Testdatei \texttt{test.php} im Hauptverzeichnis \texttt{survey} erstellt:

\lstinputlisting[
    language=php, 
    caption={Testdatei \texttt{test.php}}, 
    label={lst:test-php}
]{asset/code/Aufgabe2II.tex}

\subsection{Erwartete Ausgabe}

Bei erfolgreicher Verbindung und korrektem Datenbankinhalt zeigt die Testseite (\textit{\textsl{\url{http://localhost/survey/test.php}}}) folgende Ausgabe:

\begin{verbatim}
Verbindung erfolgreich!
Array
(
    [0] => Array
        (
            [Id] => 1
            [Name] => anonymous
        )

    [1] => Array
        (
            [Id] => 2
            [Name] => Lara
        )

    [2] => Array
        (
            [Id] => 3
            [Name] => Luisa
        )

)
\end{verbatim}

Die Ausgabe bestätigt:
\begin{itemize}
    \item Die Verbindung zur Datenbank \texttt{umfrage\_db} wurde erfolgreich hergestellt.
    \item Die Tabelle \texttt{user\_table} existiert und enthält die drei erwarteten Datensätze.
    \item Die \acs{PDO}-Abfrage (\textit{\texttt{query}}) und das Auslesen der Ergebnisse (\textit{\texttt{fetchAll}}) funktionieren korrekt.
\end{itemize}

\section{Zusammenfassung}

Die Projektstruktur wurde wie gefordert eingerichtet:

\begin{itemize}
    \item Ordner \texttt{survey} und Unterordner \texttt{\acs{inc}} wurden erstellt.
    \item Die Datei \texttt{db.php} implementiert die Datenbankverbindung mit \acs{PDO}.
    \item Die Verbindung wurde erfolgreich getestet und liefert die korrekten Daten aus der \texttt{user\_table}.
\end{itemize}

Die Datei \texttt{db.php} liegt unter folgendem Pfad:

\begin{forest}
    for tree={
        font=\ttfamily,
        grow'=0,
        child anchor=west,
        parent anchor=south,
        anchor=west,
        calign=first,
        edge path={
            \noexpand\path [draw, \forestoption{edge}] (!u.south west) +(10pt,0) |- node[fill,inner sep=1.5pt] {} (.child anchor)\forestoption{edge label};
        },
        before typesetting nodes={
            if n=1{
                insert before={[,phantom,calign with current]}
            }{}
        },
        fit=band,
        before computing xy={l=22pt},
    }
    [\textbf{C:\textbackslash xampp\textbackslash htdocs\textbackslash survey}
        [\textbf{inc\textbackslash}
            [\textbf{db.php}]
        ]
    ]
\end{forest}

Die Anbindung der Datenbank ist damit abgeschlossen. Die Anwendung kann nun auf die Datenbank zugreifen und in den folgenden Aufgaben um weitere Funktionalitäten erweitert werden.

% ==============================================================================
% =================================== Aufgabe 2 ================================
% ==============================================================================